% GENERATED FILE. Edit canonical lessons, then run npm run generate:books.
% canonical-source-hash: fnv1a64:09308cf59e0dc6d5
% canonical-lessons: HI-C346-bahas, HI-C346-tark, HI-C346-jhagarna, HI-C346-tokna, HI-C346-balki

\chapter{A Heated Argument}
\label{ch:hi-a-heated-argument}
\hlchaptermodality{346}

\begin{chapteropening}
\textbf{By the end of this chapter:} \emph{I can say who argued, and what each side said.}

Complete the last lesson of chapter 346: I can say who argued, and what each side said.
\end{chapteropening}

\section[\texorpdfstring{\hi{बहस} (bahas)}{बहस (bahas)}]{\hi{बहस} (bahas) — an argument, a debate}

\label{lesson:HI-C346-bahas}



\begin{quote}
Before the new one: say the Hindi for a means, then the Hindi for an intention.
\end{quote}

\subsection*{You'll want to know: \hi{बहस}}
\textbf{\hi{बहस}} — \emph{bahas} — \textquotedblleft{}an argument, a debate\textquotedblright{}.

\begin{quote}
\hi{कल} \hi{रात} \hi{अरुण} \hi{और} \hi{मीरा} \hi{में} \hi{लंबी} \hi{बहस} \hi{हुई।}
\end{quote}

\begin{quote}
\emph{kal rāt Aruṇ aur Mīrā meṁ lambī bahas huī.}
\end{quote}

\begin{quote}
Last night Arun and Mira had a long argument.
\end{quote}

Feminine.

\begin{cousinweb}
\textbf{\hi{बहस}} (\emph{bahas}) is usually traced to Arabic \emph{baḥth}, \textquotedblleft{}an inquiry, a discussion\textquotedblright{}, which reached Hindi through Persian.
\end{cousinweb}

\begin{grammarlens}[title={What you've built}]
One more word for giving reasons and opinions.
\end{grammarlens}

\subsection*{Your turn}
\begin{itemize}
\raggedright
  \item \emph{Say it:} \emph{bahas}
  \item \emph{Say it:} \emph{bahas}, once more
  \item \emph{Say it:} \emph{kal rāt Aruṇ aur Mīrā meṁ lambī bahas huī.} — Last night Arun and Mira had a long argument.
  \item \emph{Recall:} say the Hindi for a means, then the Hindi for an intention, then say \emph{bahas} again
\end{itemize}

\begin{tcolorbox}[breakable,colback=teal!4,colframe=teal!35!black,title={Before you move on}]
What does \hi{बहस} mean? (\textquotedblleft{}An argument, a debate\textquotedblright{}.) Say it once more.
\end{tcolorbox}
\section[\texorpdfstring{\hi{तर्क} (tark)}{तर्क (tark)}]{\hi{तर्क} (tark) — a line of reasoning, logic}

\label{lesson:HI-C346-tark}



\begin{quote}
Before the new one: say the Hindi for an intention, then the Hindi for an argument.
\end{quote}

\subsection*{You'll want to know: \hi{तर्क}}
\textbf{\hi{तर्क}} — \emph{tark} — \textquotedblleft{}a line of reasoning, logic\textquotedblright{}.

\begin{quote}
\hi{तुम्हारे} \hi{तर्क} \hi{में} \hi{एक} \hi{कमज़ोरी} \hi{है।}
\end{quote}

\begin{quote}
\emph{tumhāre tark meṁ ek kamzorī hai.}
\end{quote}

\begin{quote}
There is a weakness in your reasoning.
\end{quote}

Masculine.

\begin{cousinweb}
\textbf{\hi{तर्क}} (\emph{tark}) is a tatsama, borrowed from Sanskrit \emph{tarka}, \textquotedblleft{}reasoning, inquiry\textquotedblright{} (Monier-Williams).
\end{cousinweb}

\begin{grammarlens}[title={What you've built}]
One more word for giving reasons and opinions.
\end{grammarlens}

\subsection*{Your turn}
\begin{itemize}
\raggedright
  \item \emph{Say it:} \emph{tark}
  \item \emph{Say it:} \emph{tark}, once more
  \item \emph{Say it:} \emph{tumhāre tark meṁ ek kamzorī hai.} — There is a weakness in your reasoning.
  \item \emph{Recall:} say the Hindi for an intention, then the Hindi for an argument, then say \emph{tark} again
\end{itemize}

\begin{tcolorbox}[breakable,colback=teal!4,colframe=teal!35!black,title={Before you move on}]
What does \hi{तर्क} mean? (\textquotedblleft{}A line of reasoning, logic\textquotedblright{}.) Say it once more.
\end{tcolorbox}
\section[\texorpdfstring{\hi{झगड़ना} (jhagaṛnā)}{झगड़ना (jhagaṛnā)}]{\hi{झगड़ना} (jhagaṛnā) — to quarrel}

\label{lesson:HI-C346-jhagarna}



\begin{quote}
Before the new one: say the Hindi for an argument, then the Hindi for a line of reasoning.
\end{quote}

\subsection*{You'll want to know: \hi{झगड़ना}}
\textbf{\hi{झगड़ना}} — \emph{jhagaṛnā} — \textquotedblleft{}to quarrel\textquotedblright{}.

\begin{quote}
\hi{दोनों} \hi{भाई} \hi{खिलौने} \hi{के} \hi{लिए} \hi{झगड़ते} \hi{हैं।}
\end{quote}

\begin{quote}
\emph{donõ bhāī khilaune ke liye jhagaṛte haiṁ.}
\end{quote}

\begin{quote}
The two brothers fight over the toy.
\end{quote}

Past (he …): \hi{वह} \hi{अपने} \hi{भाई} \hi{से} \hi{झगड़ा}, \emph{vah apne bhāī se jhagṛā}.

\begin{cousinweb}
\textbf{\hi{झगड़ना}} (\emph{jhagaṛnā}) is inherited from \emph{*jhagaḍ}, \textquotedblleft{}quarrel\textquotedblright{}, a form CDIAL reconstructs from its descendants (CDIAL 5321).
\end{cousinweb}

\begin{grammarlens}[title={What you've built}]
One more word for giving reasons and opinions.
\end{grammarlens}

\subsection*{Your turn}
\begin{itemize}
\raggedright
  \item \emph{Say it:} \emph{jhagaṛnā}
  \item \emph{Say it:} \emph{jhagaṛnā}, once more
  \item \emph{Say it:} \emph{donõ bhāī khilaune ke liye jhagaṛte haiṁ.} — The two brothers fight over the toy.
  \item \emph{Recall:} say the Hindi for an argument, then the Hindi for a line of reasoning, then say \emph{jhagaṛnā} again
\end{itemize}

\begin{tcolorbox}[breakable,colback=teal!4,colframe=teal!35!black,title={Before you move on}]
What does \hi{झगड़ना} mean? (\textquotedblleft{}To quarrel\textquotedblright{}.) Say it once more.
\end{tcolorbox}
\section[\texorpdfstring{\hi{टोकना} (ṭoknā)}{टोकना (ṭoknā)}]{\hi{टोकना} (ṭoknā) — to interrupt, to cut in}

\label{lesson:HI-C346-tokna}



\begin{quote}
Before the new one: say the Hindi for a line of reasoning, then the Hindi for to quarrel.
\end{quote}

\subsection*{You'll want to know: \hi{टोकना}}
\textbf{\hi{टोकना}} — \emph{ṭoknā} — \textquotedblleft{}to interrupt, to cut in\textquotedblright{}.

\begin{quote}
\hi{जब} \hi{मीरा} \hi{बोलती} \hi{है} \hi{तब} \hi{अरुण} \hi{उसे} \hi{टोकता} \hi{है।}
\end{quote}

\begin{quote}
\emph{jab Mīrā boltī hai tab Aruṇ use ṭoktā hai.}
\end{quote}

\begin{quote}
When Mira speaks, Arun interrupts her.
\end{quote}

Past, with \hi{ने} (he or she …): \hi{उसने} \hi{मुझे} \hi{टोका}, \emph{usne mujhe ṭokā}.

\begin{cousinweb}
\textbf{\hi{टोकना}} (\emph{ṭoknā}) is inherited from \emph{*ṭōkk}, \textquotedblleft{}strike, bite\textquotedblright{}, a form CDIAL reconstructs from its descendants (CDIAL 5476). CDIAL glosses the Hindi verb \textquotedblleft{}to check, to question\textquotedblright{}.
\end{cousinweb}

\begin{grammarlens}[title={What you've built}]
One more word for giving reasons and opinions.
\end{grammarlens}

\subsection*{Your turn}
\begin{itemize}
\raggedright
  \item \emph{Say it:} \emph{ṭoknā}
  \item \emph{Say it:} \emph{ṭoknā}, once more
  \item \emph{Say it:} \emph{jab Mīrā boltī hai tab Aruṇ use ṭoktā hai.} — When Mira speaks, Arun interrupts her.
  \item \emph{Recall:} say the Hindi for a line of reasoning, then the Hindi for to quarrel, then say \emph{ṭoknā} again
\end{itemize}

\begin{tcolorbox}[breakable,colback=teal!4,colframe=teal!35!black,title={Before you move on}]
What does \hi{टोकना} mean? (\textquotedblleft{}To interrupt, to cut in\textquotedblright{}.) Say it once more.
\end{tcolorbox}
\section[\texorpdfstring{\hi{बल्कि} (balki)}{बल्कि (balki)}]{\hi{बल्कि} (balki) — but rather, on the contrary}

\label{lesson:HI-C346-balki}



\begin{quote}
Before the new one: say the Hindi for to quarrel, then the Hindi for to interrupt.
\end{quote}

\subsection*{You'll want to know: \hi{बल्कि}}
\textbf{\hi{बल्कि}} — \emph{balki} — \textquotedblleft{}but rather, on the contrary\textquotedblright{}.

\begin{quote}
\hi{यह} \hi{सस्ता} \hi{नहीं} \hi{है}, \hi{बल्कि} \hi{बहुत} \hi{महँगा} \hi{है।}
\end{quote}

\begin{quote}
\emph{yah sastā nahīṃ hai, balki bahut mahãgā hai.}
\end{quote}

\begin{quote}
It isn't cheap; on the contrary, it's very expensive.
\end{quote}

\begin{cousinweb}
\textbf{\hi{बल्कि}} (\emph{balki}) is usually traced to Persian \emph{balki}, \textquotedblleft{}but rather\textquotedblright{}.
\end{cousinweb}

\begin{grammarlens}[title={What you've built}]
One more word for giving reasons and opinions.
\end{grammarlens}

\subsection*{Your turn}
\begin{itemize}
\raggedright
  \item \emph{Say it:} \emph{balki}
  \item \emph{Say it:} \emph{balki}, once more
  \item \emph{Say it:} \emph{yah sastā nahīṃ hai, balki bahut mahãgā hai.} — It isn't cheap; on the contrary, it's very expensive.
  \item \emph{Recall:} say the Hindi for to quarrel, then the Hindi for to interrupt, then say \emph{balki} again
\end{itemize}

\begin{tcolorbox}[breakable,colback=teal!4,colframe=teal!35!black,title={Before you move on}]
What does \hi{बल्कि} mean? (\textquotedblleft{}But rather, on the contrary\textquotedblright{}.) Say it once more.
\end{tcolorbox}
